% Imporditud: tools/import_eksperimendid.py
% Allikas: Eesti füüsikaolümpiaadi eksperimendi ülesannete
%   kogu (Rael Kalda, 2023), ülesanne Ü40, lahendus L40.
\setAuthor{EFO žürii}
\setRound{piirkonnavoor}
\setYear{1999}
\setNumber{G E2}
\setDifficulty{2}
\setTopic{vedelike mehaanika}
\setVahendid{silindrilise kujuga anum, mis on poolenisti täidetud tundmatu vedelikuga, väike tühi plasttops, tuntud massiga keha, joonlaud}
\setPunktid{13}
\setAllikas{Eksperimendi kogu Ü40, lahendus lk 51}
\setLahendusseis{toores}

\prob{Vedeliku tihedus}
Määrake vedeliku tihedus.

\hint

\solu
Teooria: Vedelikku asetatud kehale mõjub üleslükkejõud F = $\rho$vgV (1 p). Keha ujub vedeliku pinnal siis, kui kehale mõjuv üleslükkejõud on arvuliselt võrdne kehale mõjuva raskusjõuga mg = $\rho$vgV (1 p), kus ruumalana on arvestatud keha vedelikus oleva osa ruumala (1 p). Ujumise tingimusest saab arvutada vedeliku tiheduse $\rho$v = mg/gV (1 p). Tähistades tühja topsiku massi $m_1$, saab tühja topsiku ujumise tingimuse kirjutada üles seosega m1g = $\rho$vgV1, kus $V_1$ on topsiku vedelikku sukeldatud osa ruumala (1 p). Kui topsikusse on asetatud tuntud massiga m$^2$ keha, vajub topsik sügavamale vedelikku ja kehtib seos ($m_1$ + m$^2$)g = $\rho$vgV2, kus $V_2$ on topsiku vedelikus oleva osa ruumala (1 p). Lahutades teisest seosest esimese ning teisendades valemit, saame vedeliku tiheduse arvutada seosest $\rho$v = m$^2$/($V_2$ $-$ $V_1$) (1 p). Koormisega ja koormiseta topsiku poolt väljatõrjutud vedeliku ruumala muudu arvutame seosest $\Delta$V = $\pi$$d_2$($l_2$ $-$ $l_1$)/4, kus kus d tähistab anuma diameetrit ning $l_1$ ja $l_2$ tähistavad vedeliku nivoosid (1 p). Kui mingi mõõdetava suuruse x suhteline mõõtmisviga on Ex = $\Delta$x/x, siis tema ruudu viga Ex2 = 2$\Delta$x/x (1 p) (korrutise suhteline viga on tegurite suhteliste vigade summa).

Mõõtmised: Asetame topsiku vedelikku ja mõõdame vedeliku nivoo. Asetame koormise topsikusse ja mõõdame uuesti vedeliku nivoo. Mõõdame anuma läbimõõdu (2 p).

Arvutused: Saadud tulemustest arvutame väljatõrjutud vedeliku ruumala muudu ning vedeliku tiheduse (1 p). Kui mõõtmist teostatakse nii, et joonlaud on vedelikus, siis tuleb arvestada ka ruumalaga, mille hõivab joonlaud.

Mõõtevea arvutamine: Vedeliku ruumala mõõtmise suhteline viga on E$\Delta$V = 2$\Delta$d/d+($\Delta$$l_2$ +$\Delta$$l_1$)/($l_2$ $-$$l_1$) (1 p).

Märkusi: Koormisena on soovitatav kasutada kuubikujulist 100 g koormist. Vedelikuks sobib vesi, millele on lisatud veidi keedusoola või värvainet, topsiks --- jogurti vms tops.
\probend
